\section{Resultados y análisis}

\subsection{Resultados de la Encuesta 1: Módulo \textit{Aprender a Aprender}}

Los resultados presentados a continuación se basan en las 38 respuestas recopiladas. Las tablas completas con los promedios por ítem y el número de respuestas se encuentran en el \hyperref[anexo:encuestas]{Anexo I} (véanse Tablas~\ref{tab:anexo-aaa-comprension} a~\ref{tab:anexo-aaa-aplicacion}).

\subsubsection*{Comprensión y uso de las herramientas del módulo}

Los datos muestran que los estudiantes comprendieron adecuadamente la estructura del módulo y el funcionamiento de sus principales herramientas. Los promedios por ítem se ubicaron en valores altos de la escala Likert (entre 4.00 y 4.39), lo que indica claridad en el uso de las metas de aprendizaje, la relación entre objetivos y tareas, el tablero Kanban, el roadmap generado por el chatbot y las autoevaluaciones (véase Tabla~\ref{tab:anexo-aaa-comprension} en el \hyperref[anexo:encuestas]{Anexo I}).

\subsubsection*{Motivación y participación}

Las actividades del módulo generaron niveles favorables de motivación y participación, con promedios que oscilaron entre 4.13 y 4.42. Elementos como el tablero Kanban, la personalización de tareas, el editor visual de roadmaps y la interacción con el chatbot fueron percibidos como factores que favorecieron el involucramiento activo durante la experiencia de uso (véase Tabla~\ref{tab:anexo-aaa-motivacion} en el \hyperref[anexo:encuestas]{Anexo I}).

\subsubsection*{Aplicación práctica y retención de la información}

Los estudiantes reportaron que el módulo apoyó la retención de conceptos clave relacionados con la planificación, la autorregulación y la organización del aprendizaje. Los promedios de esta dimensión se situaron entre 4.08 y 4.54, destacándose la afirmación sobre la aplicabilidad de los contenidos en situaciones reales como el ítem con mayor valoración. Esto sugiere que los conocimientos adquiridos pueden transferirse a contextos académicos y personales (véase Tabla~\ref{tab:anexo-aaa-aplicacion} en el \hyperref[anexo:encuestas]{Anexo I}).

\vspace{0.4cm}

\subsection{Resultados de la Encuesta 2: Módulo \textit{Escucha Activa}}

Los resultados presentados a continuación se basan en las 25 respuestas recopiladas. Las tablas completas con los promedios por ítem, resultados dicotómicos y número de respuestas se encuentran en el \hyperref[anexo:encuestas]{Anexo I} (véanse Tablas~\ref{tab:anexo-ea-comprension} a~\ref{tab:anexo-ea-global}).

\subsubsection*{Comprensión del contenido}

Los ítems asociados a la comprensión de los principios de la escucha activa mostraron valoraciones altas, especialmente en relación con la identificación de ideas principales y la aplicación de la habilidad en situaciones reales. Los promedios se ubicaron entre 4.64 y 4.80 en la escala Likert, lo que indica que las explicaciones y la retroalimentación ofrecidas por el módulo fueron percibidas como claras y útiles (véase Tabla~\ref{tab:anexo-ea-comprension} en el \hyperref[anexo:encuestas]{Anexo I}).

\subsubsection*{Motivación y participación}

Las dinámicas del módulo, que incluyen la variedad de modos de juego, la estructura por rondas, el sistema de puntajes y la posibilidad de repetir el audio, favorecieron la motivación y el involucramiento activo de los estudiantes. Los promedios oscilaron entre 4.44 y 4.72, lo que refleja que las actividades fueron percibidas como dinámicas e interesantes y que mantuvieron la atención durante las sesiones de juego (véase Tabla~\ref{tab:anexo-ea-motivacion} en el \hyperref[anexo:encuestas]{Anexo I}).

\subsubsection*{Retención de la información}

La mayoría de los estudiantes señaló haber mejorado su capacidad para recordar detalles importantes después de escuchar los audios: el 92\,\% indicó una percepción positiva de mejora y solo el 8\,\% manifestó lo contrario. Adicionalmente, la diversidad de estilos de contenido y niveles de dificultad fue valorada con un promedio de 4.52, lo que indica que, en general, los estudiantes percibieron una contribución del módulo al fortalecimiento de su comprensión auditiva. (véanse Tablas~\ref{tab:anexo-ea-retencion-dicotomica} y~\ref{tab:anexo-ea-retencion-likert} en el \hyperref[anexo:encuestas]{Anexo I}).

\subsubsection*{Aplicación práctica}

Los estudiantes reportaron niveles altos de aplicabilidad de las habilidades desarrolladas en el módulo a situaciones reales, como clases, conversaciones y trabajo en equipo. Los promedios de esta dimensión se ubicaron entre 4.44 y 4.72, lo que indica que el módulo fue percibido como útil para transferir las estrategias de escucha activa fuera del entorno de la aplicación (véase Tabla~\ref{tab:anexo-ea-aplicacion} en el \hyperref[anexo:encuestas]{Anexo I}).

\subsubsection*{Evaluación global}

La valoración global del módulo fue muy positiva. La experiencia general fue calificada con promedios superiores a 4.80, y la mayoría de los estudiantes indicó que recomendaría la aplicación a otros compañeros interesados en fortalecer su habilidad de escucha activa (véase Tabla~\ref{tab:anexo-ea-global} en el \hyperref[anexo:encuestas]{Anexo I}).

\vspace{0.4cm}

\subsection{Resultados de la Encuesta 3: Evaluación de usabilidad basada en heurísticas}

La evaluación basada en heurísticas indica que el prototipo presenta una usabilidad sólida, con valoraciones concentradas en los niveles más altos de la escala Likert (4 y 5). La Tabla~\ref{tab:usabilidad-heuristicas} sintetiza los principales hallazgos y las oportunidades de mejora identificadas.

\begin{table}[H]
\centering
\footnotesize
\setlength{\tabcolsep}{5pt}
\renewcommand{\arraystretch}{1.35}
\begin{tabular}{|p{3.2cm}|p{5.4cm}|p{5.4cm}|}
\hline
\textbf{Heurística} & \textbf{Hallazgos principales} & \textbf{Oportunidades de mejora} \\
\hline

Claridad del estado del sistema
& Valoraciones entre 4 y 5; comunicación adecuada del estado general.
& Aclarar flujos relacionados con IA y mejorar la retroalimentación en la barra lateral. \\ \hline

Lenguaje y correspondencia
& Lenguaje comprensible; sin términos confusos reportados.
& Revisar organización y denominación en la barra lateral. \\ \hline

Control y libertad
& Navegación clara; facilidad para deshacer o corregir.
& Mayor visibilidad de acciones de corrección en flujos específicos. \\ \hline

Consistencia y estándares
& Consistencia general en diseño; valoraciones altas.
& Unificar pantallas aisladas con estilos diferentes (p. ej., pantallas de éxito). \\ \hline

Prevención y manejo de errores
& Validaciones funcionales; valoraciones entre 4 y 5.
& Explicar mejor las condiciones de contraseña y mensajes de error. \\ \hline

Claridad visual y carga cognitiva
& Interfaz percibida como limpia y enfocada.
& Aumentar contraste y resaltar acciones prioritarias. \\ \hline

Ayuda y documentación
& Ayuda útil para completar tareas principales; valoraciones altas.
& Incluir guías contextuales y ejemplos, especialmente en flujos con IA. \\ \hline

\textbf{Sugerencias cualitativas generales}
& Retroalimentación positiva global; sin problemas estructurales de usabilidad.
& Mejorar visibilidad de botones clave, uniformar pantallas de éxito, fortalecer mensajes de error y ampliar guías para IA. \\ \hline

\textbf{Evaluación general}
& Usabilidad sólida percibida; interfaz coherente, clara y funcional.
& Refinar aspectos puntuales de visibilidad, claridad de mensajes y apoyo contextual. \\ \hline

\end{tabular}
\caption{Síntesis de hallazgos de la evaluación de usabilidad basada en heurísticas.}
\label{tab:usabilidad-heuristicas}
\end{table}
